\booksection{Приложение к главе 11: Профессиональные знания в области технологий изображения}

\chapterintro{Чтобы по-настоящему понять тайны мобильной фотографии и видео, нужно знать некоторые профессиональные технологии изображения. Эти знания не только помогут вам лучше использовать камеру телефона, но и дадут больше уверенности при выборе устройства.}

\booksubsection{Разрешение, частота кадров и битрейт: три составляющие качества видео}

\subsubsection*{Разрешение (Resolution):}

Разрешение определяет чёткость изображения и выражается в пикселях.

\begin{itemize}
  \item \textbf{720P (HD)}: 1280×720 пикселей, начальный уровень чёткости.
  \item \textbf{1080P (Full HD)}: 1920×1080 пикселей, самое распространённое разрешение.
  \item \textbf{2K (QHD)}: 2560×1440 пикселей, чётче 1080P.
  \item \textbf{4K (UHD)}: 3840×2160 пикселей, сверхвысокая чёткость, богатство деталей.
  \item \textbf{8K}: 7680×4320 пикселей, стандарт будущего, поддержка устройств пока ограничена.
\end{itemize}

Чем выше разрешение, тем чётче изображение, но тем больше файл и выше требования к хранилищу и вычислительной мощности.

\subsubsection*{Частота кадров (Frame Rate):}

Частота кадров определяет плавность видео, выражается в кадрах в секунду (fps).

\begin{itemize}
  \item \textbf{24fps}: стандарт кинематографа, даёт уникальное «киношное ощущение».
  \item \textbf{30fps}: обычный стандарт для ТВ и интернет-видео, плавно и естественно.
  \item \textbf{60fps}: высокая частота, подходит для спорта и записи игр, очень плавно.
  \item \textbf{120fps/240fps}: сверхвысокая частота, для замедленного видео.
\end{itemize}

Чем выше частота кадров, тем плавнее видео и лучше эффект замедления, но больше размер файла.

\subsubsection*{Битрейт (Bitrate):}

Битрейт определяет плотность данных видео, выражается в объёме данных в секунду (Мбит/с).

\begin{itemize}
  \item Чем выше битрейт, тем лучше качество и больше сохранённых деталей, но больше файл.
  \item Слишком низкий битрейт приводит к появлению артефактов сжатия — мозаики и цветовых блоков.
  \item Профессиональная видеосъёмка обычно требует высокого битрейта для сохранения пространства постобработки.
\end{itemize}

\booksubsection{Форматы видеокодеков: H.264, H.265 и ProRes}

\subsubsection*{H.264 (AVC):}

\begin{itemize}
  \item \textbf{Особенности}: самый широко используемый формат кодирования, отличная совместимость.
  \item \textbf{Плюсы}: воспроизводится почти на всех устройствах, умеренная эффективность сжатия.
  \item \textbf{Минусы}: по сравнению с новыми форматами, более низкая эффективность сжатия, файлы больше.
  \item \textbf{Сценарии}: обмен интернет-видео, повседневная съёмка обычных пользователей.
\end{itemize}

\subsubsection*{H.265 (HEVC):}

\begin{itemize}
  \item \textbf{Особенности}: обновлённая версия H.264, эффективность сжатия выше примерно на 50\%.
  \item \textbf{Плюсы}: при том же качестве файл меньше, поддержка более высоких разрешений и частот кадров.
  \item \textbf{Минусы}: кодирование и декодирование требуют более мощного процессора, некоторые старые устройства несовместимы.
  \item \textbf{Сценарии}: съёмка 4K-видео, сценарии с ограниченным хранилищем.
\end{itemize}

\subsubsection*{Apple ProRes:}

\begin{itemize}
  \item \textbf{Особенности}: профессиональный видеокодек Apple, почти без потерь.
  \item \textbf{Плюсы}: высочайшее качество, большое пространство для постобработки, точный цвет.
  \item \textbf{Минусы}: огромные файлы, требуется много места для хранения, обработка на профессиональном оборудовании.
  \item \textbf{Сценарии}: профессиональное видеопроизводство, съёмка киноуровня.
\end{itemize}

\subsubsection*{ProRes RAW:}

\begin{itemize}
  \item \textbf{Особенности}: сочетает преимущества ProRes и формата RAW.
  \item \textbf{Что такое RAW}: необработанные исходные данные сенсора, сохраняющие максимум информации.
  \item \textbf{Плюсы}: огромное пространство для постобработки — можно менять экспозицию, баланс белого и другие параметры без потери качества.
  \item \textbf{Минусы}: чрезвычайно большие файлы, требуется профессиональное ПО для обработки.
  \item \textbf{Сценарии}: профессиональное производство высокого уровня. iPhone 15 Pro и новее поддерживают съёмку Apple ProRes видео (до 4K 60fps), iPhone 17 Pro поддерживает запись ProRes RAW прямо на устройстве.
\end{itemize}

\booksubsection{HDR: технология высокого динамического диапазона}

Технология HDR (High Dynamic Range) позволяет одновременно сохранять детали в самых ярких и самых тёмных участках изображения.

\begin{itemize}
  \item \textbf{HDR10}: базовый стандарт HDR, хорошая совместимость.
  \item \textbf{HDR10+}: стандарт, продвигаемый Samsung, поддержка динамических метаданных.
  \item \textbf{Dolby Vision}: премиальный стандарт HDR от Dolby, лучшее качество, но требует лицензирования.
  \item \textbf{HLG}: вещательный стандарт HDR, совместим с обычными мониторами.
\end{itemize}

Для просмотра полного эффекта HDR-видео требуется HDR-совместимый дисплей.

\booksubsection{Цветовое пространство и глубина цвета}

\subsubsection*{Цветовое пространство:}

\begin{itemize}
  \item \textbf{sRGB}: стандартное цветовое пространство, наилучшая совместимость.
  \item \textbf{Display P3}: продвигаемый Apple стандарт расширенного цветового охвата, более насыщенные цвета.
  \item \textbf{Rec.2020}: стандартное цветовое пространство для HDR-видео, самый широкий охват.
\end{itemize}

\subsubsection*{Глубина цвета:}

\begin{itemize}
  \item \textbf{8-bit}: 16,7 млн цветов, стандарт обычного видео.
  \item \textbf{10-bit}: 1,07 млрд цветов, стандарт профессионального видео, более плавные цветовые переходы.
  \item \textbf{12-bit}: 68 млрд цветов, профессиональный стандарт высокого уровня.
\end{itemize}

\begin{tcbinfo}{Практические рекомендации}
\begin{itemize}
  \item Повседневная съёмка: 1080P/4K + 30fps + кодирование H.265 — баланс качества и размера файла.
  \item Съёмка движения: 1080P/4K + 60fps — для плавности.
  \item Профессиональное творчество: 4K + ProRes (если поддерживается) — максимум пространства для постобработки.
  \item Совет по хранению: видео с высоким битрейтом занимает много места, рекомендуется регулярно создавать резервные копии на компьютер или в облако.
\end{itemize}

\end{tcbinfo}

\begin{tcbpractice}{Исследуйте: настройки видео на вашем телефоне}
\begin{enumerate}
  \item Откройте приложение «Камера», переключитесь в режим видео.
  \item Проверьте настройки разрешения и частоты кадров, попробуйте разные комбинации.
  \item Если поддерживается, проверьте наличие HDR, ProRes и других продвинутых опций.
  \item Снимите несколько видео с разными настройками и сравните качество и размер файлов.
\end{enumerate}

\end{tcbpractice}

\booksubsection{Сравнение форматов RAW, HEIF и JPEG}

\subsubsection*{Что такое формат RAW:}

Формат RAW — это «исходные данные», записанные сенсором камеры, — чистая необработанная информация без сжатия и обработки. Это как сырые продукты, требующие «приготовления» при постобработке. Из-за ограниченного хранилища и вычислительной мощности телефона по умолчанию снимки сохраняются уже «приготовленными» в JPEG или HEIF, RAW нужно включать специально.

\subsubsection*{Сравнение RAW, HEIF и JPEG:}

\begin{itemize}
  \item \textbf{JPEG}: старейший формат, лучшая совместимость (открывается на любом устройстве). 8-битная глубина цвета, сжатие с потерями, самый маленький файл (2-5 МБ на фото). Ограниченное пространство для постобработки — при сильном осветлении теней появляются цветовые блоки.
  \item \textbf{HEIF/HEIC}: новый формат, продвигаемый Apple, при том же качестве файл вдвое меньше JPEG (1-3 МБ на фото). Поддержка 10-битной глубины цвета и альфа-канала. Недостаток: некоторые старые устройства и Windows требуют установки плагинов для открытия.
  \item \textbf{RAW}: несжатые исходные данные, самый большой файл (20-100 МБ на фото). Глубина цвета обычно 12-16 бит, максимальное пространство для постобработки. Недостаток: обязательно профессиональное ПО (Lightroom, Snapseed и т.д.) для открытия и обработки, нельзя напрямую поделиться.
\end{itemize}

\subsubsection*{Почему профессиональные фотографы снимают в RAW:}

\begin{itemize}
  \item \textbf{Баланс белого можно без потерь скорректировать при постобработке}: если JPEG снят с неправильным балансом белого (например, желтизна при лампах накаливания), исправление приведёт к потере качества; RAW позволяет перенастроить цветовую температуру полностью без потерь.
  \item \textbf{Большой допуск по экспозиции}: в JPEG недодержка на 2 ступени в тенях — практически мёртвая чернота, после осветления — сплошной шум; RAW может «спасти» полезные детали из теней при недодержке 3-4 ступени. То же с передержкой.
  \item \textbf{Богаче детали в тенях и светах}: динамический диапазон RAW значительно шире JPEG — больше запаса, чтобы света не были пересвечены, а тени не были мёртво-чёрными.
  \item \textbf{Резкость, шумоподавление, коррекция объектива — всё контролируется при постобработке}: внутрикамерная обработка — это «единоразовая сделка», RAW можно настраивать многократно, каждый раз отталкиваясь от исходных данных.
\end{itemize}

\subsubsection*{Как включить RAW на телефоне:}

\begin{itemize}
  \item \textbf{iPhone ProRAW}: только для iPhone 12 Pro и новее Pro-моделей. Включите в «Настройки Камера Форматы» опцию «Apple ProRAW». Затем в интерфейсе камеры сверху появится кнопка «RAW» — нажмите для активации. Файл около 25 МБ/снимок, рекомендуется включать только при необходимости.
  \item \textbf{Android Motion Raw / Ultra RAW}: названия различаются у брендов. Флагманы Xiaomi/OPPO/vivo — переключатель RAW в «Профессиональном режиме» или «Ещё»; Huawei — включение RAW в «Профессиональном режиме»; Samsung Galaxy S — опция RAW в «Профессиональном» режиме. Размер файла обычно 20-50 МБ.
  \item \textbf{Общий совет}: для повседневных случайных снимков RAW не нужен — файлы слишком большие; включайте только для важных сцен (пейзаж, портрет, требуется постобработка).
\end{itemize}

\booksubsection{Принципы вычислительной фотографии}

Современная мобильная фотография уже неотделима от «вычислительной фотографии» — то есть создания снимка с помощью алгоритмов, а не только объектива и сенсора. Понимание этих ключевых принципов поможет вам лучше осознать, что делает телефон в момент нажатия кнопки спуска.

\subsubsection*{Многокадровое объединение:}

При одном нажатии затвора телефон на самом деле делает несколько снимков. Алгоритм выравнивает кадры, взвешенно объединяет их и получает одно высококачественное изображение.

\begin{itemize}
  \item \textbf{Ночной режим}: за 1-3 секунды делается серия недоэкспонированных кадров, они выравниваются и объединяются — шумы взаимно гасятся, яркость накапливается. Результат — чистая картинка даже при слабом освещении.
  \item \textbf{HDR (высокий динамический диапазон)}: снимаются недо-, нормально и переэкспонированные кадры (три или более), затем объединяются в один снимок с деталями и в светах, и в тенях. Особенно эффективен при контровом свете.
  \item \textbf{Deep Fusion (Apple) / Photonic Engine (фотонный движок)}: для сцен со средним и низким освещением — попиксельный анализ нескольких кадров и их слияние с акцентом на повышение текстуры и детализации. Детали вроде шерсти и волос становятся заметно чётче.
\end{itemize}

\subsubsection*{Семантическая сегментация:}

AI не просто обрабатывает всё изображение «одним махом», а сначала «понимает», что на нём изображено — небо, лицо, трава, здание, одежда — и затем применяет разные стратегии обработки к каждой области.

\begin{itemize}
  \item \textbf{Небо}: отдельно усиливается синева, снижается яркость во избежание пересвета.
  \item \textbf{Лицо}: распознаются участки кожи, отдельно осветляется и сглаживается тон кожи, не затрагивая черты лица и волосы.
  \item \textbf{Трава/листва}: усиливается насыщенность зелёного, повышается резкость деталей.
  \item \textbf{Размытие фона}: в портретном режиме AI распознаёт контур объекта — фон размыт, объект чёток.
\end{itemize}

\subsubsection*{Применение глубокого обучения:}

\begin{itemize}
  \item \textbf{AI-распознавание сцен}: камера автоматически определяет сцены — «еда», «животное», «ночь», «документ» — и соответственно оптимизирует параметры. Например, режим «Еда» усиливает тёплые тона и насыщенность, делая блюдо более аппетитным.
  \item \textbf{AI-размытие портрета}: объектив телефона физически не может дать малую глубину резкости как зеркальная камера — модель глубокого обучения предсказывает границы объекта и размывает фон. Новейшие алгоритмы уже корректно обрабатывают такие детали, как пряди волос и края стекла.
  \item \textbf{AI-удаление / AI-расширение}: выделите прохожего или лишний предмет в кадре — AI на основе контекста «дорисует» скрытую область. Это применение генеративного AI в области изображений.
\end{itemize}

\subsubsection*{Вычислительная фотография vs оптическая фотография:}

\begin{itemize}
  \item \textbf{Преимущества вычислительной фотографии}: хорошее качество при малом размере, возможность съёмки в экстремальных условиях (слабое освещение, контровой свет), обычный пользователь без знаний фотографии получает «красивые» снимки.
  \item \textbf{Недостатки вычислительной фотографии}: иногда «переобработка» — избыточная резкость, искажение цвета, неестественный цвет неба; многокадровое объединение оставляет ореолы при движении объекта; детали «вычислены», а не записаны реально.
  \item \textbf{Преимущества оптической фотографии}: большой сенсор + светосильный объектив реально захватывают свет, естественное качество, большое пространство для постобработки; отсутствие артефактов от «домысливания» алгоритмов.
  \item \textbf{Тенденция}: современные флагманские телефоны сочетают «оптику + вычисления» — большой сенсор получает реальную информацию, алгоритмы оптимизируют представление. Это не противостояние, а синергия.
\end{itemize}

\booksubsection{Типы объективов и фокусные расстояния}

Современные флагманские телефоны обычно оснащены тремя-четырьмя объективами, каждый под своё фокусное расстояние. Понимание фокусных расстояний поможет выбрать правильный объектив для разных сцен.

\subsubsection*{Распространённые объективы и эквивалентное фокусное расстояние:}

\begin{itemize}
  \item \textbf{Сверхширокоугольный (13-24 мм экв.)}: самый широкий угол обзора, подходит для пейзажей, архитектуры, интерьеров. Захватывает масштабную сцену, но по краям есть искажения (прямые линии искривляются), при портрете лицо вытягивается. Для «широких сцен», не для «крупных планов».
  \item \textbf{Основной (24-28 мм экв.)}: самый используемый объектив, лучшее качество, максимальная светосила. Баланс угла обзора и качества — и для людей, и для пейзажей. «1x» на большинстве телефонов — это основной объектив.
  \item \textbf{Телеобъектив (70-200 мм экв.)}: для удалённых объектов и портретных крупных планов. Телеобъектив сжимает пространство (далёкая гора и близкий человек выглядят ближе друг к другу), даёт более естественные пропорции лица и размытие фона.
  \item \textbf{Перископический телеобъектив (200+ мм экв.)}: за счёт отражающей оптической схемы достигает большего фокусного расстояния в пределах толщины телефона. Позволяет чётко снимать детали с большого расстояния — концерты, спорт, дикая природа. Только во флагманских моделях.
\end{itemize}

\subsubsection*{Компромисс между размером сенсора и размером пикселя:}

Сенсор — это «основа», определяющая, сколько света может быть принято. В индустрии говорят: «больший сенсор подавляет всё» — при одинаковом количестве пикселей чем больше сенсор, тем больше отдельный пиксель, тем больше света, тем лучше качество.

\begin{itemize}
  \item \textbf{Большой сенсор}: много света, хорошее качество ночных снимков, широкий динамический диапазон. Недостаток: телефон становится толще (выступающий модуль камеры), дороже.
  \item \textbf{Высокое разрешение на маленьком сенсоре}: при дневном свете высокое разрешение даёт более чёткие снимки, удобно для последующего кадрирования. Но при слабом свете каждому пикселю не хватает света — много шума.
  \item \textbf{Выбор флагманов}: основной сенсор обычно 1/1,3 дюйма до 1 дюйма (некоторые модели уже используют 1-дюймовый сенсор), разрешение около 50 Мп — это точка баланса качества и размера.
\end{itemize}

\subsubsection*{Технология объединения пикселей (Pixel Binning):}

Чтобы совместить высокое разрешение днём и большие пиксели ночью, производители используют «объединение пикселей» — несколько маленьких пикселей сливаются в один большой для вывода.

\begin{itemize}
  \item \textbf{4-в-1 (Quad Bayer)}: сенсор 50 Мп, вывод 12,5 Мп. Каждый выходной пиксель формируется из 4 соседних пикселей, световой поток увеличен в 4 раза, качество чище.
  \item \textbf{9-в-1 (Nonacell)}: для сенсоров выше 100 Мп, вывод около 11-12 Мп. Тот же принцип, ещё больше света.
  \item \textbf{Реальный эффект}: в обычном режиме съёмки вы получаете объединённое фото с низким разрешением — хорошее качество, маленький файл. Когда нужна предельная детализация (например, съёмка документа или дальнего плана), можно вручную переключиться в «Режим высокого разрешения» для полного вывода.
\end{itemize}

\booksubsection{Специальные режимы съёмки}

Помимо повседневной съёмки, в телефоне есть несколько специальных режимов — при правильном использовании можно снять то, что обычно не получается.

\subsubsection*{Панорамный режим:}

\begin{itemize}
  \item При съёмке медленно и плавно перемещайте телефон горизонтально по стрелке-подсказке, не спешите — при быстром движении стыковка будет со смещением.
  \item Держите телефон стабильно по вертикали, двигайте только горизонтально. Если качать вверх-вниз, линия горизонта станет волнистой.
  \item Осторожно при наличии движущихся людей или машин — при склейке могут появиться «полтела» или «призраки».
  \item Панораму можно снимать и вертикально — вертикальная съёмка с горизонтальной склейкой подходит для высоких зданий и водопадов.
\end{itemize}

\subsubsection*{Замедленная и интервальная съёмка:}

\begin{itemize}
  \item \textbf{Замедленная съёмка}: съёмка на 120fps или 240fps, при воспроизведении — замедленное видео. Подходит для капель воды, спортивных моментов, прыжков животных. Требуется достаточное освещение — в режиме slow-motion время экспозиции короткое, при слабом свете качество плохое.
  \item \textbf{Интервальная съёмка (timelapse)}: съёмка серии кадров с интервалом и объединение в ускоренное видео. Подходит для автомобильных потоков, движения облаков, заката, распускания цветов. Длительность от нескольких минут до нескольких часов — обязательно используйте штатив, с рук гарантированно смазано.
  \item \textbf{Совет}: при интервальной съёмке облаков установите телефон, снимайте 15-30 минут — получится около 10 секунд timelapse, эффект впечатляющий.
\end{itemize}

\subsubsection*{Звёздное небо / Астро-режим:}

\begin{itemize}
  \item Нужно уехать подальше от городского светового загрязнения, найти место без уличных фонарей — за городом или в горах. В городе Млечный Путь не снять.
  \item Обязательно штатив, с рук гарантированно смазано. Телефон должен быть неподвижен от нескольких секунд до десятков секунд.
  \item Погода ясная, без облаков, луна не слишком большая (полнолуние заглушает звёздный свет). Лучше всего — тёмные ночи в начале или конце лунного месяца.
  \item «Ночной режим» iPhone автоматически распознает звёздное небо и увеличивает выдержку (до 30 сек.); у флагманов Android аналогичный эффект в «Режиме звёздного неба» или «Астро-режиме». Наберитесь терпения — после нажатия спуска не трогайте телефон.
\end{itemize}

\subsubsection*{Макро-режим:}

\begin{itemize}
  \item Для съёмки цветов, насекомых, ювелирных изделий, текста крупным планом и других мелких объектов. Можно фокусироваться с расстояния 1-2 см.
  \item Некоторые телефоны имеют отдельный макро-объектив (обычно 2 Мп или 5 Мп на маленьком сенсоре), качество среднее; флагманы используют сверхширокоугольный объектив для макро — качество лучше.
  \item При съёмке вплотную объектив находится близко к объекту и загораживает свет — при необходимости используйте настольную лампу или фонарик для подсветки.
  \item Фокусное расстояние очень маленькое, глубина резкости минимальна — смещение объекта на несколько миллиметров вперёд-назад может сбить фокус. Рука должна быть твёрдой, лучше чтобы объект оставался неподвижным.
\end{itemize}

\begin{tcbinfo}{Практические советы по статической фотографии}
\begin{itemize}
  \item Повседневная съёмка: HEIF/JPEG по умолчанию — маленький файл, достаточное качество.
  \item Важные сцены (пейзаж, портрет, требуется постобработка): включите RAW — максимум пространства для постобработки.
  \item Правильный выбор объектива: пейзаж — сверхширик, портрет — телеобъектив, повседневность — основной.
  \item Ночь и звёзды: обязательно зафиксируйте телефон, включите ночной или астро-режим.
  \item Телефоны с большим сенсором (например, 1-дюймовые модели) имеют более высокий потолок качества, но привычки оператора остаются ключевым фактором.
\end{itemize}

\end{tcbinfo}

\begin{tcbpractice}{Попробуйте сами: изучите возможности вашей камеры}
\begin{enumerate}
  \item Проверьте, поддерживает ли ваш телефон съёмку в RAW (Настройки Камера Формат), попробуйте включить и сделать снимок, сравните размер файлов JPEG и RAW для одной сцены.
  \item Переключайтесь между сверхширокоугольным, основным и телеобъективом, снимая один и тот же дальний план — почувствуйте разницу в перспективе и качестве.
  \item Найдите сцену для ночного или астро-режима, зафиксируйте телефон и сделайте ночной снимок.
  \item Попробуйте снять панораму в горизонтальной ориентации — следите за тем, чтобы телефон двигался плавно движущимся.
\end{enumerate}

\end{tcbpractice}
